166
12 Surfaces and Differential Forms in $\mathbb{R}^n$
In this case the mapping
$$
\begin{array}{l}
t^1 = f^1(x^1,\ldots,x^k), \\
\vdots \\
t^k = f^k(x^1,\ldots,x^k), \\
t^{k+1} = x^{k+1} - f^{k+1}(x^1,\ldots,x^k), \\
\vdots \\
t^n = x^n - f^n(x^1,\ldots,x^k)
\end{array}
$$
is a diffeomorphism of a full-dimensional neighborhood of the point $x_0 \in \mathbb{R}^n$. As $\varphi_{\varepsilon}$
we can now take the restriction to some cube $I^n$ of the diffeomorphism inverse to
it.
By a change of scale, of course, one can arrange to have $\varepsilon = 1$ and a unit cube
$I^n$ in the last diffeomorphism.
Thus we have shown that for a smooth surface in $\mathbb{R}^n$ one can adopt the following
definition, which is equivalent to the previous one.

Definition 4 The $k$-dimensional surface in $\mathbb{R}^n$ introduced by Definition 1 is smooth
(of class $C^{(m)}, m \geq 1$) if it has an atlas whose local charts are smooth mappings (of
class $C^{(m)}, m \geq 1$) and have rank $k$ at each point of their domains of definition.

We remark that the condition on the rank of the mappings (12.1) is essential. For
example, the analytic mapping $\mathbb{R} \ni t \mapsto (x^1, x^2) \in \mathbb{R}^2$ defined by $x^1 = t^2, x^2 = t^3$
defines a curve in the plane $\mathbb{R}^2$ having a cusp at $(0, 0)$. It is clear that this curve is
not a smooth one-dimensional surface in $\mathbb{R}^2$, since the latter must have a tangent
(a one-dimensional tangent plane) at each point.$^4$

Thus, in particular one should not conflate the concept of a smooth path of class
$C^{(m)}$ and the concept of a smooth curve of class $C^{(m)}$.
In analysis, as a rule, we deal with rather smooth parametrizations (12.1) of
rank $k$. We have verified that in this case Definition 4 adopted here for a smooth
surface agrees with the one considered earlier in Sect. 8.7. However, while the pre-
vious definition was intuitive and eliminated certain unnecessary complications im-
mediately, the well-known advantage of Definition 4 of a surface, in accordance
with Definition 1, is that it can easily be extended to the definition of an abstract
manifold, not necessarily embedded in $\mathbb{R}^n$. For the time being, however, we shall be
interested only in surfaces in $\mathbb{R}^n$.
Let us consider some examples of such surfaces.
$^4$For the tangent plane see Sect. 8.7.